\chapter{Analysis of Results}
\label{chap:analysis}

The measurements presented in \cref{sec:results-presentation} show clear differences between configurations on some metrics and largely overlapping distributions on others. Despite this discrepancy, several consistent patterns emerge across the benchmarks and metrics, which allow for a structured interpretation of the results.

\section{Non-Strong Reference Processing}
\label{sec:analysis-non-strong}

The metric most directly targeted by the implemented optimisations is \texttt{Concurrent Process Non-Strong}, which captures the total wall-clock time spent in ZGC's concurrent reference-processing phase per GC cycle. Because this phase is the intended site of all three optimisations, the clearest differences between variants are observed here, making it the natural starting point for the analysis.

\Cref{fig:gc-single-nonstrong,fig:gc-multi-nonstrong} show the distributions of this metric for the single-object and multi-object benchmarks respectively. Full numerical summaries are given in \cref{tab:stats-nonstrong}.

\subsection{Effect of the Clear Path and Dynamic Array}
\label{subsec:analysis-clear-path-dyn}

The most striking reduction is achieved by the \texttt{clear\_path\_dyn} variant. In the single-object benchmark its median non-strong processing time is 187.9\,ms, compared to 996.9\,ms for the unmodified \texttt{none} baseline, a reduction of approximately 81\%. The \texttt{all} variant, which adds the skip-enqueue separation on top of \texttt{clear\_path\_dyn}, reaches a nearly identical median of 184.4\,ms (also about an 81\% reduction). In the multi-object benchmark the pattern is consistent: \texttt{clear\_path\_dyn} achieves a median of 18.8\,ms versus 44.1\,ms for \texttt{none} (57\% reduction), and \texttt{all} reaches 18.7\,ms (also 57\% reduction).

The near-equivalence of \texttt{clear\_path\_dyn} and \texttt{all} is a notable finding. Adding the skip-enqueue separation to a configuration that already has the dynamic array and the specialised clear path yields little extra benefit: \texttt{sep\_only} reduces the single-object median by only 5\%, while \texttt{clear\_path\_dyn} and \texttt{all} differ by just 3.5\,ms in the single-object benchmark and 0.1\,ms in the multi-object benchmark.

The two individual components tell a different story. The \texttt{clear\_path\_only} variant, which enables only the specialised clear path, reduces the single-object median by a modest 7\% (to 923.7\,ms), while \texttt{dyn\_only}, which enables only the dynamic array, achieves a 36\% reduction (to 639.1\,ms). When combined in \texttt{clear\_path\_dyn}, the effect is substantially larger than either isolated contribution. This suggests that the dynamic-array traversal and the specialised clear logic remove different parts of the per-reference cost and therefore interact constructively rather than redundantly.

\subsection{Effect of the Skip-Enqueue Separation Alone}
\label{subsec:analysis-sep-only}

The \texttt{sep\_only} variant, which introduces the separate discovered list for queue-less weak references but does not include the dynamic array or the specialised clear path, produces negligible improvement on non-strong processing time. In the single-object benchmark its median is 943.8\,ms, only 5\% below the \texttt{none} baseline, and its IQR of 73.8\,ms is comparable to \texttt{none}'s 75.9\,ms. In the multi-object benchmark the difference is entirely negligible: 44.1\,ms versus 44.1\,ms for \texttt{none}, well within measurement variability.

\subsection{Comparison with Weak Fields}
\label{subsec:analysis-non-strong-weak-fields}

The \texttt{weak\_fields} variant achieves a median non-strong processing time of 484.6\,ms in the single-object benchmark and 35.8\,ms in the multi-object benchmark. These figures correspond to reductions of 51\% and 19\% respectively relative to \texttt{none}, positioning \texttt{weak\_fields} clearly above the \texttt{none} baseline but well below the leading \texttt{clear\_path\_dyn} and \texttt{all} configurations. Its non-strong phase therefore improves clearly over the baseline, but not enough to match the best specialised \texttt{WeakReference} pipeline variants. This contrast is important because \texttt{weak\_fields} still performs very strongly on broader GC metrics (\cref{sec:analysis-inter-model}), which indicates that its main advantage lies outside the non-strong subphase itself.

\section{Overall GC Collection Times}
\label{sec:analysis-overall-gc}

In contrast to the non-strong processing metric, the overall collection time metrics, major collection time and old-generation collection time, show broad, heavily overlapping distributions across all \texttt{WeakReference} variants, making it difficult to draw firm conclusions from pairwise comparisons within this group. \Cref{fig:gc-single-major,fig:gc-multi-major} illustrate this for the major collection metric. Numerical summaries for all variants are given in \cref{tab:stats-major,tab:stats-oldgen,tab:stats-markfollow}.

\subsection{WeakReference Variants}
\label{subsec:analysis-weakref-overall}

For the single-object benchmark, the median major collection time of \texttt{none} is 6\,958\,ms with an interquartile range (IQR) of about 1\,596\,ms. The \texttt{all} variant has a median of 6\,377\,ms (8\% below \texttt{none}) and the \texttt{clear\_path\_dyn} variant 6\,399\,ms (also 8\% below \texttt{none}), but with IQRs of 1\,419\,ms and 1\,394\,ms respectively, both distributions still overlap substantially with the baseline and with each other. Several variants that do not combine the strongest non-strong optimisations, such as \texttt{sep\_only} and \texttt{clear\_path\_only}, show medians at or above \texttt{none}. In the multi-object benchmark the effect is even smaller in relative terms, with medians ranging from 965\,ms (\texttt{dyn\_only}) to 1\,002\,ms (\texttt{clear\_path\_sep}) against a baseline of 982\,ms; all distributions remain within a narrow IQR band of roughly 93--136\,ms, making fine-grained ranking unreliable.

Among the \texttt{WeakReference} variants, \texttt{all} and \texttt{clear\_path\_dyn} tend to occupy the lower part of the collection-time distribution in both benchmarks, which is directionally consistent with their superior non-strong processing times. However, given the distribution overlap, this observation should be treated as indicative rather than conclusive.

\subsection{Violin Plot Shape and the Bimodal Pattern}
\label{subsec:analysis-bimodal}

Many violin plots for the overall collection metrics exhibit a distinctly bimodal shape, with two separated density peaks rather than a single bell-shaped distribution. This pattern is visible in both benchmarks but is more pronounced for the single-object case, where per-run maxima are used. The distribution separation is not caused by any particular variant: all nine configurations exhibit it. The same broad two-regime shape is also visible in the mark-follow tables, whose single-object medians cluster around 2.6--2.8 seconds for the \texttt{WeakReference} variants while their IQRs remain very large (560--810\,ms). This makes it plausible that variability in concurrent marking work, not just the non-strong phase, is driving much of the spread in total collection time.

\subsection{Single versus Multi-Object Benchmark}
\label{subsec:analysis-single-vs-multi}

The single-object and multi-object benchmarks yield broadly consistent rankings across the key non-strong metric, giving confidence that the strongest effects are not artefacts of one particular workload. The principal difference is in variability: the single-object benchmark produces substantially taller violin plots and a higher incidence of outliers. For end-to-end collection metrics the ordering is less stable, but the same broad picture remains: the best \texttt{WeakReference} variants offer only modest total-time gains, whereas \texttt{weak\_fields} shifts the entire GC cost profile more substantially.

\section{Inter-Model Comparison: Weak Fields versus WeakReference}
\label{sec:analysis-inter-model}

The \texttt{weak\_fields} variant stands apart from all \texttt{WeakReference} variants on the overall GC collection metrics (see \cref{tab:stats-major,tab:stats-oldgen,tab:stats-markfollow}). In the single-object benchmark, its median major collection time is 4\,136\,ms, compared to 6\,377--7\,044\,ms for the \texttt{WeakReference} variants, a reduction of approximately 35--41\% depending on which variant it is compared against. For old-generation collection time the pattern is similar: 2\,295\,ms versus 2\,993--3\,676\,ms, corresponding to a reduction of about 23--38\%. In the multi-object benchmark the gap is smaller but still clear: 708\,ms versus 965--1\,002\,ms for major collection (a reduction of about 27--29\%), and 554\,ms versus 792--831\,ms for old generation (about 29--33\%). The \texttt{weak\_fields} distributions are also notably tighter, with an IQR of 485\,ms for single-object major collection, compared with 1\,354--1\,596\,ms for the \texttt{WeakReference} variants. The concurrent mark-follow phase shows a similar pattern: a median of 1\,782\,ms for \texttt{weak\_fields} versus 2\,575\,ms for \texttt{none} in the single-object benchmark, a reduction of 31\%.

Taken together, these numbers indicate that the principal advantage of weak fields is structural rather than procedural. The weak-field mechanism does not produce the lowest non-strong processing time, but it removes a large population of wrapper objects from the heap and thereby reduces the amount of GC work outside the reference-processing subphase as well.

\section{Memory Behaviour}
\label{sec:analysis-memory}

\subsection{Auxiliary Memory Usage and the Dynamic Array}
\label{subsec:analysis-aux-memory}

\Cref{fig:mem-single-hist,fig:mem-multi-hist} show the median percentage difference in per-run maximum GC-committed (auxiliary) memory relative to the \texttt{none} baseline across all configurations. Absolute values are listed in \cref{tab:stats-memory-aux}.

All configurations that include the dynamic array (\texttt{dyn\_only}, \texttt{sep\_dyn}, \texttt{clear\_path\_dyn}, and \texttt{all}) incur substantially higher auxiliary memory than variants that retain the intrusive linked list. In the single-object benchmark, the median per-run maximum for \texttt{dyn\_only} and \texttt{sep\_dyn} is 308\,MB of GC-committed memory, compared to approximately 120\,MB for the linked-list variants. The \texttt{clear\_path\_dyn} variant is the highest of all at 1\,268\,MB, and \texttt{all} reaches 884\,MB. In the multi-object benchmark the absolute differences are smaller, with medians of 292--363\,MB across all variants.

The \texttt{weak\_fields} variant occupies an intermediate position. Its median auxiliary memory in the single-object benchmark is 446\,MB, higher than the non-array \texttt{WeakReference} variants (approximately 120\,MB) but lower than \texttt{clear\_path\_dyn}. In the multi-object benchmark it is the lowest of all configurations at 240\,MB.

\subsection{Java Heap Usage}
\label{subsec:analysis-heap}

The Java heap usage panel in \cref{fig:mem-single-hist,fig:mem-multi-hist} shows the median percentage difference in per-run maximum heap usage relative to the \texttt{none} baseline. Absolute values are listed in \cref{tab:stats-memory-java-heap}.

The heap usage figures are essentially identical across all \texttt{WeakReference} variants (approximately 1\,720\,MB per-run median in the single-object benchmark and 1\,926--1\,928\,MB in the multi-object benchmark).

The \texttt{weak\_fields} variant, however, shows markedly lower heap usage: a median of approximately 806\,MB in the single-object benchmark (roughly half of the \texttt{WeakReference} variants) and 1\,834\,MB in the multi-object benchmark.

The RSS results follow the same overall direction as auxiliary memory and heap usage. In the single-object benchmark, \texttt{clear\_path\_dyn} reaches the highest median RSS at 2\,856\,MB, compared with 1\,901\,MB for \texttt{none}, whereas \texttt{weak\_fields} reaches only 1\,260\,MB. In the multi-object benchmark, the spread is narrower, from 2\,255\,MB for \texttt{weak\_fields} to 2\,487\,MB for \texttt{clear\_path\_dyn}. This strengthens the interpretation that the dynamic-array variants trade runtime memory overhead for faster queue-less weak-reference processing, while weak fields reduce both heap occupancy and total process footprint.
